\section{Estudos de Caso}
\label{sec:case_studies}

Para demonstrar a viabilidade da metodologia proposta, o modelo foi aplicado em duas capitais brasileiras: Curitiba (PR) e Porto Alegre (RS). A escolha desses municípios permite testar os dois comportamentos da função de penalização definidos na Seção \ref{sec:penalization}: o impacto de perturbações ambientais contínuas (temperatura de superfície) e discretas (áreas de risco geológico).

Para garantir a reprodutibilidade da pesquisa, os experimentos priorizam a utilização de dados abertos e bases governamentais oficiais.

\subsection{Curitiba: Degradação por Variável Contínua (Clima)}
O experimento em Curitiba avaliará como o desconforto térmico afeta a mobilidade ativa. Neste cenário, a temperatura de superfície atuará como um penalizador contínuo do esforço de deslocamento, aumentando o peso das arestas mais expostas ao calor. Os dados utilizados compreendem:

\begin{itemize}
    \item \textbf{Malha de deslocamento:} Rede viária extraída do OpenStreetMap (OSM), filtrada para infraestrutura caminhável.
    \item \textbf{Pontos de Interesse (POIs):} Coletados via OSM, complementados e validados por bases de alvarás municipais da Prefeitura de Curitiba, categorizados por tipos de serviço.
    \item \textbf{Condicionante ambiental:} Imagens de satélite (\textit{rasters}) do produto \textit{Surface Temperature} obtidas via Earth Explorer (Landsat 8 e 9, sensor TIRS - banda ST\_TRAD).
\end{itemize}

\subsection{Porto Alegre: Degradação por Variável Discreta (Risco Físico)}
O experimento em Porto Alegre analisará o impacto de áreas de risco geotécnico e físico sobre a acessibilidade urbana. Como condicionantes ambientais, serão utilizadas camadas vetoriais (polígonos) de \textit{Áreas com Risco a Movimentos de Massa} e \textit{Áreas de Risco Geológico}, ambas disponíveis no portal GeoPOA\footnote{\url{https://gis-smamus.portoalegre.rs.gov.br/dmweb-2-0/}} e resultantes de mapeamento realizado pelo Serviço Geológico do Brasil (SCG/CPRM). Essas áreas atuam como variáveis discretas, simulando barreiras que interrompem ou penalizam rotas afetadas por riscos. Os dados considerados para o experimento são:

\begin{itemize}
    \item \textbf{Malha de deslocamento:} Rede viária extraída do OpenStreetMap (OSM);
    \item \textbf{Pontos de Interesse (POIs):} Serviços e instalações mapeados integralmente a partir do OSM;
    \item \textbf{Condicionantes ambientais:} Polígonos das \textit{Áreas com Risco a Movimentos de Massa} e \textit{Áreas de Risco Geológico} proveniente do GeoPOA (origem: SCG/CPRM).
\end{itemize}

Por fim, para a análise de desigualdade socioeconômica em ambos os municípios, as malhas territoriais serão enriquecidas com dados do Censo Demográfico do IBGE (Instituto Brasileiro de Geografia e Estatística), permitindo o cruzamento das perdas de acessibilidade com variáveis de renda, infraestrutura domiciliar e vulnerabilidade social.
